% ============================================
% The Opus Audio Codec: Theory and Implementation
% Educational manual for RFC 6716 and the opus_native pure-Rust implementation
% ============================================

\section{Introduction}\label{sec:introduction}

This manual provides a comprehensive overview of the algorithms,
optimizations, and design decisions that make up the Opus audio
codec~\cite{rfc6716}. The focus is on understanding both the theoretical
foundations and the practical implementations of the techniques Opus combines:
the information theory of its entropy coder, the linear prediction of its speech
path, the transform coding of its music path, and the fixed-point and SIMD
machinery that lets all of it run in real time.

The material emphasizes the \termdef{why} and \termdef{how} of each component:
\emph{why} certain design choices matter, and \emph{how} they are implemented
efficiently. It covers both the mathematics and the practical computational
considerations that make these methods work in production systems.

Opus is the IETF's general-purpose, royalty-free audio codec, standardized in
2012 and updated by RFC 8251~\cite{rfc8251}. It is the default audio codec of
WebRTC and is widely deployed in streaming and conferencing. Its defining
characteristic is that it is not one codec but two, combined behind a single
bitstream. A linear-predictive speech coder (\termdef{SILK}) and a
transform-domain music coder (\termdef{CELT}, the Constrained-Energy Lapped
Transform) share one entropy coder, one packet format, and one container. The
encoder selects, per frame, between SILK, CELT, or both together (the
\termdef{hybrid} mode). The combination spans a range no single technique covers
well: 6~kb/s narrowband speech to 510~kb/s fullband stereo music, frame sizes
from 2.5~ms to 60~ms, and sample rates from 8~kHz to 48~kHz, all decoded by one
receiver that is not told in advance which mode each packet uses.

\subsection{The Reference Implementation}\label{subsec:reference_impl}

Throughout, the implementation under discussion is
\texttt{opus\_native}~--- a pure-Rust, dependency-free realisation of RFC 6716.
Every code listing in this manual is drawn from that source. Three of its design
choices shape the discussion:

\begin{itemize}[itemsep=3pt]
  \item \textbf{Bit-exact by construction.} The bitstream is defined by the
    range coder (\cref{sec:rangecoder}), not by any one encoder's choices. The
    implementation's entropy arithmetic follows the RFC byte-for-byte, and the
    decoder is validated against the official test vectors on the
    \emph{final-range oracle} --- a per-packet checksum of the entropy-coder
    state that must agree with the reference to the bit.
  \item \textbf{Safe by default.} The crate denies \code{unsafe} everywhere
    except a handful of explicitly annotated SIMD kernels
    (\cref{sec:optimizations}). The signal-processing math is ordinary safe
    Rust; the only machine intrinsics sit in encoder hot loops where a
    faster, approximate result is provably harmless (\cref{sec:optimizations}).
  \item \textbf{Decode-first.} The decoder is the normative half of the
    specification and the half the conformance vectors exercise. The encoder is
    free to make any choice the decoder can follow; this asymmetry recurs
    throughout the manual and is the key to several optimisations.
\end{itemize}

\subsection{Target Audience and Prerequisites}\label{subsec:prerequisites}

This material assumes familiarity with:
\begin{itemize}[itemsep=3pt]
  \item Discrete signals, sampling, and the discrete Fourier transform
  \item Linear algebra (vectors, matrices, inner products, projection)
  \item Basic probability and the idea of entropy / information content
  \item Programming and algorithm analysis; the listings are in Rust, but no
    Rust expertise is required to follow them
\end{itemize}

No prior knowledge of speech coding, transform coding, arithmetic coding, or
psychoacoustics is assumed; each is introduced where it is first needed. A
reader who has worked through a companion treatment of FFT-based spectral
analysis will find the transform-coding sections (CELT) immediately familiar,
and the linear-prediction sections (SILK) a complementary story told in the
time domain.

\subsection{Scope and Organisation}\label{subsec:organization}

The manual is organised bottom-up, from the bit to the file, mirroring the
layering of the codec itself:

\begin{enumerate}[itemsep=3pt]
  \item \textbf{Architecture Overview} (\cref{sec:architecture}) --- the two
    codecs, three modes, the shared range coder, and the sample-rate model.
  \item \textbf{The Range Coder} (\cref{sec:rangecoder}) --- entropy coding,
    arithmetic/range coding, the ICDF representation, raw bits, \code{tell()},
    and the final-range oracle.
  \item \textbf{Packet Framing} (\cref{sec:packet}) --- the TOC byte, the 32
    configurations, frame packing codes 0--3, padding, and validity rules.
  \item \textbf{CELT: Transform Coding} (\cref{sec:celt}) --- the MDCT, the
    Vorbis window and TDAC, pyramid vector quantisation, CWRS enumeration, band
    splitting, energy coding, bit allocation, and the encoder's analysis.
  \item \textbf{SILK: Linear-Predictive Coding} (\cref{sec:silk}) --- linear
    prediction, LSF/NLSF quantisation, long-term (pitch) prediction, the
    noise-shaping quantiser, gains, the shell coder, and resampling.
  \item \textbf{The Opus Layer} (\cref{sec:opus_layer}) --- mode dispatch,
    the SILK/CELT crossover, redundancy at transitions, packet-loss
    concealment, in-band FEC, and discontinuous transmission.
  \item \textbf{Ogg Encapsulation} (\cref{sec:ogg}) --- pages, the Ogg CRC,
    lacing, packet reassembly, \code{OpusHead}/\code{OpusTags}, and granule
    timing.
  \item \textbf{Multistream and Surround} (\cref{sec:multistream}) --- coupled
    and uncoupled streams, channel mapping families.
  \item \textbf{Computational Optimisations} (\cref{sec:optimizations}) ---
    SIMD dot products, the SIMD PVQ search, the MDCT pre-rotation, the FFT
    backend seam, table-driven CWRS, and buffer reuse.
  \item \textbf{Numerical Considerations and Conformance}
    (\cref{sec:numerical}) --- fixed-point versus float, bit-exactness, the
    oracle, and the conformance criterion.
\end{enumerate}

\subsection{Notation Conventions}\label{subsec:notation}

Throughout this manual:
\begin{itemize}[itemsep=2pt]
  \item $x[n]$ denotes a time-domain sample at index $n$; $X[k]$ a
    frequency-domain coefficient at bin $k$
  \item $f_s$ denotes a sample rate in Hz; Opus's internal reference rate is
    always 48~kHz
  \item $N$ denotes a transform length or vector dimension; $K$ a pulse count
  \item $\mathrm{LM}$ (``log multiplier'') denotes $\log_2$ of the frame-size
    multiplier: $\mathrm{LM}=0,1,2,3$ for 2.5, 5, 10, 20~ms CELT frames
  \item The range coder's state is the pair $(\mathit{val}, \mathit{rng})$
  \item \termdef{\Q{n}} denotes a fixed-point value scaled by $2^n$: a \Q{15}
    value $v$ represents the real number $v / 2^{15}$
  \item $\sgn(x)$ is the sign of $x$; $\floor*{\cdot}$ and $\ceil*{\cdot}$ are
    floor and ceiling
  \item Section references like ``\S4.3.4'' without further qualification refer
    to RFC 6716~\cite{rfc6716}
\end{itemize}

\section{Architecture Overview}\label{sec:architecture}

This section describes the codec as a whole --- its components, the path a
packet takes through them, and the design decisions that the rest of the manual
elaborates --- before later sections treat each layer in detail.

\subsection{Two Codecs, One Bitstream}\label{subsec:two_codecs}

Opus is built on the observation that audio divides into two regimes that reward
opposite coding strategies. The codec implements both and selects between them
per frame.

\begin{concept}{Speech versus Music: Two Coding Strategies}
  \textbf{Speech} is produced by a physical system --- the vocal folds exciting
  the vocal tract --- that is well modelled as a slowly varying \emph{linear
  filter} driven by a sparse \emph{excitation}. Coding speech efficiently means
  coding that model: estimate the filter, estimate the pitch, and transmit only
  the residual that the model fails to predict. This is \termdef{linear-predictive
  coding} (LPC), and it is the approach SILK takes, in the \emph{time domain}, at
  internal rates of 8, 12, or 16~kHz.

  \textbf{Music} is arbitrary: dense, polyphonic, and percussive, with energy
  spread across the full 20~kHz audible band and no simple source model. Coding
  music efficiently means working in the \emph{frequency domain}, where the
  ear's own frequency analysis allows bits to be spent where they are audible
  and withheld where masking hides them. This is \termdef{transform coding}, and
  it is the approach CELT takes, via the Modified Discrete Cosine Transform
  (MDCT) at 48~kHz.

  Neither technique performs both tasks well: LPC degrades on polyphonic
  material, and transform coding spends bits inefficiently on the strong
  predictable structure of speech. Opus therefore retains both and switches
  between them per frame; in the overlap region (wideband speech with a musical
  background) it runs them together as the \termdef{hybrid} mode.
\end{concept}

This yields the three \termdef{modes}, selected by the packet header
(\cref{sec:packet}) and dispatched by the decoder (\cref{sec:opus_layer}):

\begin{center}
  \begin{tabular}{lll}
    \toprule
    \textbf{Mode} & \textbf{Coder(s)} & \textbf{Typical use} \\
    \midrule
    SILK-only & SILK & low-bitrate speech, 8--16~kHz audio band \\
    Hybrid    & SILK (low) + CELT (high) & wideband speech + full bandwidth \\
    CELT-only & CELT & music, low delay, high bitrate \\
    \bottomrule
  \end{tabular}
\end{center}

In hybrid mode the split is by frequency: SILK codes the signal up to roughly
8~kHz (it runs at a 16~kHz internal rate), and CELT codes the band above,
starting at MDCT band 17 (\cref{subsec:hybrid}). Both coders write into the same
range coder, so there is no boundary marker between them in the byte stream: the
CELT data is simply a continuation of the same range-coded message that began
with the SILK data.

\subsection{The Layered Stack}\label{subsec:stack}

\Cref{fig:stack-table} lists the layers as the implementation modularises them,
bottom to top. The lower layers (range coder, packet, Ogg) are pure framing and
entropy --- exact integer arithmetic, \code{no\_std}-compatible, with no audio
math. The middle layers (SILK, CELT) perform the signal processing. The top
layer (the Opus decoder and encoder) provides the glue: mode dispatch,
resampling, and the cross-frame state used to handle mode transitions.

\begin{table}[h]
  \centering
  \begin{tabular}{lll}
    \toprule
    \textbf{Layer} & \textbf{Specification} & \textbf{Responsibility} \\
    \midrule
    \code{ogg}         & RFC 3533 / 7845 & file container, pages, timing \\
    \code{multistream} & RFC 7845 \S5.1.1 & surround: $N$ coupled/mono streams \\
    \code{decoder}/\code{encoder} & RFC 6716 \S4/\S5 & mode dispatch, transitions, PLC/FEC/DTX \\
    \code{silk}        & RFC 6716 \S4.2 & linear-predictive speech coder \\
    \code{celt}        & RFC 6716 \S4.3 & MDCT transform coder \\
    \code{packet}      & RFC 6716 \S3 & TOC byte, frame packing \\
    \code{range}       & RFC 6716 \S4.1/\S5.1 & entropy coder \\
    \bottomrule
  \end{tabular}
  \caption{The Opus layer stack, bottom to top.}
  \label{fig:stack-table}
\end{table}

\subsection{Decoding a Packet: An Overview}\label{subsec:packet_life}

To decode one Opus packet, the decoder performs the following steps:

\begin{enumerate}[itemsep=3pt]
  \item Reads the first byte --- the \termdef{TOC} (Table of Contents) --- which
    names the mode, audio bandwidth, frame size, and channel count
    (\cref{sec:packet}).
  \item Splits the remaining bytes into 1--48 \emph{frames} according to the
    TOC's frame-packing code.
  \item For each frame, constructs a range decoder over its bytes and, depending
    on the mode, runs the SILK decoder, the CELT decoder, or both into a shared
    range decoder.
  \item Resamples the decoder's internal output (always referenced to 48~kHz) to
    the caller's requested rate, and applies any cross-frame smoothing for mode
    transitions and redundancy (\cref{sec:opus_layer}).
\end{enumerate}

The top-level dispatch is a compact \code{match} on the parsed TOC:

\begin{lstlisting}[language=Rust]
pub fn decode_packet(&mut self, data: &[u8]) -> Result<Vec<f32>, PacketError> {
    if data.len() <= 1 {
        // Empty or TOC-only packet: conceal one frame (DTX / loss).
        return Ok(self.decode_lost(self.last_frame_size / self.ds));
    }
    let packet = Packet::parse(data)?;
    let toc = packet.toc();
    let mut out = Vec::new();
    for frame in packet.frames() {
        let pcm = self.decode_frame(
            frame, toc.mode(), toc.bandwidth(),
            toc.frame_size().samples_per_channel_48k(),
        );
        out.extend_from_slice(&pcm);
    }
    Ok(out)
}
\end{lstlisting}

\subsection{Sample Rates and the 48 kHz Reference}\label{subsec:internal_rate}

The relationship between sample rates in Opus is a common source of confusion.
The model is stated once here and referred to throughout.

\begin{concept}{The 48 kHz Reference Clock}
  Opus has three distinct notions of sample rate:
  \begin{itemize}[itemsep=2pt]
    \item The \textbf{audio bandwidth} (narrowband, mediumband, wideband,
      super-wideband, fullband) is a \emph{property of the content}: how much of
      the spectrum carries signal. It is named in the TOC and caps the highest
      coded frequency (4, 6, 8, 12, or 20~kHz).
    \item The \textbf{internal rate} is the rate a coder actually runs at. SILK
      runs at 8, 12, or 16~kHz; CELT always runs at 48~kHz.
    \item The \textbf{API rate} is the rate the caller asks for --- 8, 12, 16,
      24, or 48~kHz --- at which PCM is delivered.
  \end{itemize}

  Crucially, \textbf{all timing is expressed in 48~kHz samples} regardless of
  the other two. A 20~ms frame is always 960 samples ``at 48~kHz''; a
  2.5~ms frame is 120. Granule positions in the Ogg container
  (\cref{sec:ogg}) count 48~kHz samples. The decoder resamples from the internal
  rate up to 48~kHz (CELT is already there), then \emph{down} to the API rate by
  an integer factor $\mathit{ds} = 48000 / f_s$. This single reference clock is
  what lets SILK frames, CELT frames, and the container agree on ``when'' a
  sample occurs even though they compute it at different rates.
\end{concept}

\subsection{From Quantities to Symbols}\label{subsec:bits_overview}

Every numeric quantity Opus transmits --- a gain, an LSF index, a pulse vector,
an energy delta --- becomes a symbol in the range coder. The encoder and decoder
share static probability models (\termdef{ICDF} tables, \cref{sec:rangecoder})
for the structured fields and fall back to \termdef{raw bits} for the
near-uniform ones. This arrangement has an important consequence:

\begin{importantnote}
  The bitstream is defined by the \emph{sequence of range-coder symbols}, not
  by the floating-point intermediate values the encoder computed to choose them.
  Two encoders that pick the same symbols produce identical bytes, even if one
  used SIMD and \code{rsqrt} approximations and the other used scalar
  double-precision. This is why the decoder must be bit-exact (it is normative)
  but the encoder need only be \emph{valid} (any choice the decoder can follow
  is legal). Most of \cref{sec:optimizations} exploits exactly this freedom.
\end{importantnote}

\section{The Range Coder}\label{sec:rangecoder}

Every bit Opus transmits passes through the range coder. Because the bitstream
is defined entirely by the sequence of range-coder symbols, and because
bit-exact conformance is verified through the coder's internal state
(\cref{subsec:oracle}), it is the natural starting point and is treated here
first and in detail.

\subsection{Why Entropy Coding}\label{subsec:why_entropy}

Suppose a symbol $s$ occurs with probability $p(s)$. Information theory says the
\emph{ideal} cost of transmitting it is its \termdef{self-information},
$-\log_2 p(s)$ bits~\cite{shannon1948mathematical,cover2006elements}. A common
symbol ($p$ near 1) should cost a fraction of a bit; a rare one, many bits. Over
a long message drawn from a known distribution, the average cost approaches the
\termdef{entropy} $H = -\sum_s p(s)\log_2 p(s)$, and no lossless code can do
better.

A fixed-length code cannot reach this bound: it must spend a whole number of
bits per symbol, wasting the fractional part. \termdef{Arithmetic coding}, and
its integer-friendly cousin \termdef{range coding}~\cite{martin1979range,
witten1987arithmetic}, do reach it, by abandoning the idea that each symbol maps
to its own bit pattern. Instead, the \emph{entire message} is encoded as a
single number in $[0, 1)$, and each symbol narrows the interval that number must
lie in.

\begin{concept}{Range Coding by Interval Subdivision}
  Begin with the interval $[0, 1)$. To encode a symbol $s$ whose cumulative
  distribution carves out the sub-range $[\mathit{fl}/\mathit{ft},\,
  \mathit{fh}/\mathit{ft})$ (low and high cumulative frequencies out of a total
  $\mathit{ft}$), \emph{replace} the current interval with that fraction of
  itself:
  \begin{equation}\label{eq:rc_subdivide}
    [\,\ell, \ell + r\,) \;\longmapsto\;
    \left[\,\ell + r\tfrac{\mathit{fl}}{\mathit{ft}},\;
            \ell + r\tfrac{\mathit{fh}}{\mathit{ft}}\,\right).
  \end{equation}
  The interval shrinks by a factor of $(\mathit{fh}-\mathit{fl})/\mathit{ft} =
  p(s)$ at each step, so after the whole message its width is $\prod_i p(s_i) =
  2^{-\sum_i -\log_2 p(s_i)}$. Transmitting \emph{any} number inside the final
  interval takes about $-\log_2(\text{width}) = \sum_i -\log_2 p(s_i)$ bits ---
  the entropy. The decoder, knowing the same model, replays the subdivisions:
  it sees which sub-range the received number falls in, emits that symbol,
  rescales, and repeats.
\end{concept}

The practical obstacle is that the interval bounds need ever more precision as
they shrink, which is impossible with finite-width integers. The fix, due to
Pasco and refined by Martin, is \termdef{renormalisation}: as soon as the
interval is narrow enough that its top bits are settled, \emph{emit those bits}
and shift them out, restoring headroom. Opus renormalises a byte at a time.

\subsection{The Integer Range Coder}\label{subsec:integer_rc}

Opus's range coder (\S4.1 for the decoder, \S5.1 for the encoder) works in
32-bit integers. Its state is two words:

\begin{itemize}[itemsep=2pt]
  \item \code{rng} --- the current interval width, kept in
    $(2^{23}, 2^{31}]$.
  \item \code{val} --- on the encoder, the interval's low end; on the decoder,
    the offset of the received number from the top of the interval.
\end{itemize}

The governing constants name the renormalisation threshold and the byte width:

\begin{lstlisting}[language=Rust]
const SYM_BITS:  u32 = 8;            // emit/consume one byte at a time
const SYM_MAX:   u32 = 255;          // 2^8 - 1
const CODE_BITS: u32 = 32;           // width of the (val, rng) state
const CODE_TOP:  u32 = 1 << 31;      // 2^31: initial encoder range
const CODE_BOT:  u32 = 1 << 23;      // 2^23: renormalise when rng falls to here
const CODE_SHIFT: u32 = 23;          // CODE_BITS - SYM_BITS - 1; carry bit position

#[inline]
const fn ilog(x: u32) -> u32 { u32::BITS - x.leading_zeros() }  // 1 + floor(log2 x)
\end{lstlisting}

The lower bound $\mathit{rng} > 2^{23}$ guarantees that the division
$\mathit{rng}/\mathit{ft}$ in \cref{eq:rc_subdivide} keeps enough significant
bits for the largest total frequency Opus uses ($\mathit{ft} \le 2^{16}$). Note
\code{ilog} is just the hardware count-leading-zeros instruction; it appears
everywhere the coder needs $\floor*{\log_2}$.

\subsubsection{Decoding a symbol}

Decoding is two steps: \code{decode} reports \emph{where} the received value
sits in $[0, \mathit{ft})$, the caller maps that to a symbol, then \code{update}
shrinks the state to that symbol's sub-range and renormalises.

\begin{lstlisting}[language=Rust]
pub fn decode(&mut self, ft: u32) -> u32 {
    // Locate val within [0, ft): the symbol s satisfies fl <= decode() < fh.
    ft - (self.val / (self.rng / ft) + 1).min(ft)
}

pub fn update(&mut self, fl: u32, fh: u32, ft: u32) {
    let s = self.rng / ft;
    self.val -= s * (ft - fh);            // shift past the symbols above this one
    self.rng = if fl > 0 { s * (fh - fl) } // interior symbol: exact sub-width
               else       { self.rng - s * (ft - fh) }; // first symbol keeps remainder
    self.normalize();
}
\end{lstlisting}

Two details are easy to miss and important to get right. First, the first symbol
($\mathit{fl}=0$) is given the \emph{remainder} of the range
($\mathit{rng} - s\,(\mathit{ft}-\mathit{fh})$), not $s\cdot\mathit{fh}$; this is
how the coder accounts for the truncation in $s = \floor*{\mathit{rng}/
\mathit{ft}}$ and stays reversible. Second, \code{decode} caps its result at
$\mathit{ft}$ via \code{.min(ft)} to defend against a corrupt stream pushing the
quotient out of range.

\subsubsection{The renormalisation loop}

When \code{rng} drops to $2^{23}$ or below, the decoder shifts in a fresh byte:

\begin{lstlisting}[language=Rust]
fn normalize(&mut self) {
    while self.rng <= CODE_BOT {
        self.nbits_total += SYM_BITS;
        self.rng <<= SYM_BITS;
        // Reassemble the next 8-bit symbol from the leftover bit of the
        // previous byte (high) and the top 7 bits of the new byte (low).
        let byte = self.read_byte();
        let sym = (self.leftover_bit << 7) | (byte >> 1);
        self.leftover_bit = byte & 1;
        self.val = ((self.val << SYM_BITS) + (SYM_MAX - sym)) & 0x7FFF_FFFF;
    }
}
\end{lstlisting}

The one-bit shimmy (\code{leftover\_bit}) is a peculiarity of Opus's coder: the
decoder works with a 7-bits-from-this-byte, 1-bit-carried representation so that
the encoder's carry propagation (below) lands on byte boundaries. The
\code{\& 0x7FFF\_FFFF} keeps \code{val} within 31 bits.

\subsubsection{Encoding a symbol}

The encoder mirrors the decoder exactly --- it must, or the two will disagree
and the stream will be undecodable:

\begin{lstlisting}[language=Rust]
pub fn encode(&mut self, fl: u32, fh: u32, ft: u32) {
    let r = self.rng / ft;
    if fl > 0 {
        self.val += self.rng - r * (ft - fl);
        self.rng  = r * (fh - fl);
    } else {
        self.rng -= r * (ft - fh);
    }
    self.normalize();
}
\end{lstlisting}

The encoder's \code{normalize} emits the top byte of \code{val} instead of
reading one, but is otherwise the same loop --- with one complication that
deserves its own treatment.

\subsection{Carry Propagation}\label{subsec:carry}

When the encoder adds to \code{val} (the $\mathit{fl}>0$ branch above), the
addition can overflow into bits it has \emph{already emitted}. A byte written
out as \code{0xFF} might need to become \code{0x00} with a carry rippling into
the byte before it. The encoder cannot un-write bytes, so it does not write them
eagerly. Instead it buffers:

\begin{lstlisting}[language=Rust]
fn carry_out(&mut self, c: u32) {
    if c == SYM_MAX {
        self.ext += 1;                 // a run of 0xFF: defer, might all carry
    } else {
        let carry = (c >> SYM_BITS) as u8;
        if let Some(rem) = self.rem.take() {
            self.write_byte(rem + carry);     // resolve the buffered byte
        }
        if self.ext > 0 {
            // Flush the deferred 0xFF run, +carry turns them into 0x00.
            let sym = (SYM_MAX + u32::from(carry)) as u8;
            for _ in 0..self.ext { self.write_byte(sym); }
            self.ext = 0;
        }
        self.rem = Some((c & SYM_MAX) as u8);
    }
}
\end{lstlisting}

\begin{concept}{Why a 0xFF Run Must Be Held}
  A carry into a \code{0xFF} byte turns it into \code{0x00} and \emph{propagates
  further left}. A carry into any other byte simply increments it and stops. So
  the encoder can safely emit a non-\code{0xFF} byte only once it knows whether a
  carry is coming --- which it learns from the \emph{next} byte. It therefore
  keeps one pending byte in \code{rem} and counts a run of \code{0xFF}s in
  \code{ext}. When a non-\code{0xFF} byte arrives with carry $c$: the pending
  byte becomes $\mathit{rem}+c$, and the held \code{0xFF}s become either
  \code{0xFF} (no carry) or \code{0x00} (carry). This buffering is invisible to
  the decoder; it simply reads the bytes that finally emerge. The mechanism is
  what makes the encoder/decoder \emph{symmetric to the bit}, which is the
  property the conformance oracle (\cref{subsec:oracle}) checks.
\end{concept}

\subsection{The ICDF Representation}\label{subsec:icdf}

Structured fields are coded against static probability tables. Opus stores them
not as a PDF or a forward CDF but as an \termdef{inverse cumulative distribution
function} (ICDF): for a model with total $2^{\mathit{ftb}}$, the table holds
$\mathit{icdf}[k] = 2^{\mathit{ftb}} - \mathit{fh}[k]$ --- the cumulative
\emph{tail} above each symbol, decreasing to a terminating zero.

\begin{concept}{Why Inverse CDFs}
  Two reasons. First, the decode primitive in \cref{eq:rc_subdivide} naturally
  produces ``how far down from the top is \code{val}'', which is exactly an
  upper-tail comparison --- so an ICDF lets the decoder walk the table with one
  multiply and one compare per step and no running subtraction. Second, the
  tables are tiny (a handful of bytes), and storing the tail lets the common
  ``most probable symbol first'' ordering terminate early. The forward
  arithmetic ($\mathit{fl}, \mathit{fh}$) is recovered implicitly from adjacent
  ICDF entries.
\end{concept}

\begin{lstlisting}[language=Rust]
pub fn decode_icdf(&mut self, icdf: &[u8], ftb: u32) -> usize {
    let d = self.val;
    let r = self.rng >> ftb;          // ft is a power of two: shift, don't divide
    let mut k = 0usize;
    let mut t = self.rng;
    let mut s = r * u32::from(icdf[0]);
    while d < s {                     // walk down the tail until val sits in [s, t)
        k += 1;
        t = s;
        s = r * u32::from(icdf[k]);
    }
    self.val = d - s;                 // fold update() into the same step
    self.rng = t - s;
    self.normalize();
    k
}
\end{lstlisting}

Because $\mathit{ft} = 2^{\mathit{ftb}}$ is a power of two, the division
$\mathit{rng}/\mathit{ft}$ becomes a shift --- a small but pervasive saving,
since ICDF decoding is the hottest path in SILK and in CELT's side information.

\subsection{Binary Symbols and Raw Bits}\label{subsec:binary_raw}

Two specialisations round out the coder.

\textbf{Binary symbols with a power-of-two probability.} Many flags have a
probability of the form $2^{-\mathit{logp}}$. These skip the table entirely:

\begin{lstlisting}[language=Rust]
pub fn decode_bit_logp(&mut self, logp: u32) -> bool {
    let s = self.rng >> logp;        // sub-range for the "1" outcome
    let bit = self.val < s;
    if !bit { self.val -= s; }
    self.rng = if bit { s } else { self.rng - s };
    self.normalize();
    bit
}
\end{lstlisting}

\textbf{Raw bits.} Some quantities are effectively uniform --- a sign, the low
bits of a pulse split --- and modelling them as probabilities wastes effort.
Opus codes these as \termdef{raw bits}, written at the \emph{opposite end} of
the same buffer from the arithmetic-coded data.

\begin{concept}{The Two-Ended Buffer}
  A range-coded packet has \emph{two} write heads in one byte array. The
  arithmetic coder fills from the front, most-significant first, byte by byte as
  it renormalises. Raw bits fill from the \emph{back}, least-significant first.
  The two heads grow toward each other; the packet is well-formed as long as
  they do not cross. This costs nothing --- raw bits need no probability model,
  so they need not interleave with the arithmetic stream --- and it lets the
  finaliser (\cref{subsec:tell}) pack the last arithmetic byte and the last raw
  byte into the same physical byte when they happen to fit. The decoder reads
  raw bits backwards from the end of the packet through a small 32-bit window,
  refilling a byte at a time.
\end{concept}

\begin{lstlisting}[language=Rust]
pub fn decode_raw_bits(&mut self, bits: u32) -> u32 {
    if self.nend_bits < bits {                       // refill the window from the end
        loop {
            self.end_window |= self.read_byte_from_end() << self.nend_bits;
            self.nend_bits += SYM_BITS;
            if self.nend_bits > WINDOW_SIZE - SYM_BITS { break; }
        }
    }
    let ret = self.end_window & ((1 << bits) - 1);    // LSB-first
    self.end_window >>= bits;
    self.nend_bits  -= bits;
    ret
}
\end{lstlisting}

\textbf{Uniform integers} combine the two: a value in $[0, \mathit{ft})$ with
$\mathit{ft}$ not a power of two is split into an arithmetic-coded high part
(8 bits of model) and a raw-bit low part.

\begin{lstlisting}[language=Rust]
pub fn encode_uint(&mut self, t: u32, ft: u32) {
    let ftb = ilog(ft - 1);
    if ftb <= 8 {
        self.encode(t, t + 1, ft);                    // small: pure arithmetic
    } else {
        let ft_hi = ((ft - 1) >> (ftb - 8)) + 1;
        self.encode(t >> (ftb - 8), (t >> (ftb - 8)) + 1, ft_hi);
        self.encode_raw_bits(t & ((1 << (ftb - 8)) - 1), ftb - 8);  // low bits raw
    }
}
\end{lstlisting}

\subsection{Bit Budgeting: tell() and tell\_frac()}\label{subsec:tell}

Opus's bit allocation (\cref{subsec:celt_allocation}) needs to know, mid-frame,
exactly how many bits have been spent so far --- and it needs that figure to
\emph{agree} between encoder and decoder, since both run the same allocation
arithmetic from the same running total. The coder exposes this through
\code{tell} and its eighth-of-a-bit refinement \code{tell\_frac}.

\begin{lstlisting}[language=Rust]
pub fn tell(&self) -> u32 { self.nbits_total - ilog(self.rng) }

pub(crate) fn tell_frac(nbits_total: u32, rng: u32) -> u32 {
    let mut lg = ilog(rng);
    let mut r = rng >> (lg - 16);     // Q15 fractional part of log2(rng)
    for _ in 0..3 {                   // 3 iterations -> 1/8-bit resolution
        r = (r * r) >> 15;
        let bit = r >> 16;
        lg = 2 * lg + bit;
        r >>= bit;
    }
    nbits_total * 8 - lg
}
\end{lstlisting}

\begin{concept}{Bits Used: Total Shifted In Minus Range Headroom}
  The intuition: \code{nbits\_total} counts every bit shifted into the coder.
  But the coder has not \emph{committed} all of them --- the current range
  \code{rng} still spans $\log_2(\mathit{rng})$ bits of undecided value. So the
  bits actually consumed are \code{nbits\_total} $- \log_2(\mathit{rng})$.
  \code{tell()} uses the integer $\floor*{\log_2}$ (a conservative whole-bit
  count); \code{tell\_frac()} refines $\log_2(\mathit{rng})$ to 3 fractional bits
  by the Newton-style squaring loop above, giving the $1/8$-bit resolution Opus's
  allocator works in. The two readouts agree on encoder and decoder because they
  are computed from the same $(\mathit{nbits\_total}, \mathit{rng})$.
\end{concept}

\subsection{The Final-Range Oracle}\label{subsec:oracle}

The range coder provides a single, powerful correctness check for any Opus
implementation, described below.

\begin{concept}{rng-Agreement and the Final Range}
  After encoding (or decoding) the same sequence of symbols, the encoder and
  decoder must hold \emph{identical} \code{rng} values --- and, transitively,
  identical \code{val} and identical output bytes. RFC 6716 elevates this to a
  conformance tool: \code{OPUS\_GET\_FINAL\_RANGE} returns the coder's \code{rng}
  at the end of each packet, and a correct decoder's final range must equal the
  reference encoder's, \emph{exactly}, for every packet.

  This is a remarkably strong check. The entire decode --- thousands of symbols
  across SILK and CELT, every table lookup, every renormalisation --- collapses
  into a single 32-bit number that is wrong if \emph{any} step diverged. The
  \texttt{opus\_native} conformance suite enforces it on all twelve official
  vectors: a passing decode is not merely close to the reference, it is
  arithmetically identical through the entropy coder. The PCM may differ by a
  rounding bit in the final float MDCT; the \emph{range} may not differ at all.
\end{concept}

In hybrid and redundancy cases (\cref{sec:opus_layer}) the packet carries more
than one range-coded segment; the implementation XORs their final ranges
together to form the single reported value, matching the reference's bookkeeping.

% ===END-OF-MANUAL===
